With increasingly frequent cross-regional mobility and sustained growth in road freight demand in China, expressway operation and management are shifting from merely ensuring traffic flow to refined governance oriented toward safety, efficiency, and resilience. In expressway scenarios, vehicles travel at high speeds, traffic flow is dense, and disturbances propagate rapidly, which creates an urgent need for continuous perception and refined management capabilities on expressway sections. The development of Intelligent Transportation Systems (ITS) has made real-time traffic-state perception and service support possible, and trajectory data that continuously describe the spatiotemporal motion of vehicles have become an important basis for traffic analysis, prediction, and control. Therefore, how to transform the discrete detection results produced by roadside sensors into continuous and reliable vehicle trajectories has become a key issue to be addressed in expressway vehicle target tracking.

Geomagnetic sensors are suitable for large-scale, low-cost deployment because of their low cost, small size, flexible deployment, and low sensitivity to illumination and weather conditions. However, two types of problems still arise in practical applications. One is measurement disorder, short-term missed detections, and false alarms caused by communication link fluctuations, noise interference, adjacent-lane interference, and parallel driving of two vehicles, which reduce online tracking accuracy. The other is consecutive missed detections caused by the degradation of sensor detection performance after long-term continuous operation, which easily lead to trajectory fragmentation and affect trajectory continuity and identity consistency. To address these problems, this thesis investigates a dual-lane vehicle target tracking method based on multiple geomagnetic sensors by deploying geomagnetic sensors on both sides of a dual-lane roadway to collect vehicle information. The main work of this thesis is as follows.

To address measurement arrival disorder, short-term missed detections, noise interference, and complex parallel-vehicle scenarios in a dual-lane multi-geomagnetic-sensor vehicle target tracking system, an online multi-vehicle target tracking method based on geomagnetic sensor state adaptation and probabilistic scenario discrimination is proposed. First, a time-window buffering strategy is adopted to correct the arrival order of measurements, and a fixed-lag window is used to admit late-arriving measurements. Combined with a local replay mechanism, this strategy improves the completeness of trajectory information. Second, based on the vehicle motion model and the geomagnetic measurement model, a vehicle state estimation method is constructed within a recursive Kalman filtering framework to realize online prediction and updating of vehicle position and velocity. Then, the tracking gate and association likelihood are modeled, and association weights are calculated by combining the posterior estimates of geomagnetic sensor detection probability and clutter intensity with geomagnetic disturbance peak consistency information, thereby achieving robust association between vehicle trajectories and multi-geomagnetic-sensor measurements in cluttered environments. Finally, a three-state scenario discrimination model for adjacent-lane interference and parallel driving of two vehicles is constructed, and the driving lane is identified and corrected online in combination with recursive vehicle state updating. Experimental results show that the proposed method can achieve relatively accurate estimation of vehicle position, velocity, and driving lane under disordered, missing, and noisy geomagnetic sensor data, and can effectively distinguish between the two easily confused scenarios of parallel two-vehicle driving and adjacent-lane interference.

To address the long-gap trajectory fragmentation caused by the degraded detection capability of some geomagnetic sensors and the resulting increase in consecutive missed detections after long-term continuous system operation, a real-time vehicle trajectory clustering algorithm based on multiple geomagnetic sensors for consecutive missed-detection scenarios is further proposed. Taking the trajectory fragments produced by upstream online tracking and the unassociated single-point measurements as inputs, the method constructs a cluster graph suitable for online incremental processing, and forms a fragment-level trajectory reconstruction process around cluster initialization, clustering similarity scoring, online cluster assignment, cluster state updating, and lifecycle maintenance. Experimental results show that the proposed method can effectively recover broken trajectory fragments under consecutive missed-detection conditions, reduce erroneous stitching and identity switches, and improve trajectory continuity and identity consistency.